\section{Lecture 23: Attention Mechanisms}

\subsection{Motivation}

In the previous lecture, we studied sequence models based on recurrence.

\medskip

\noindent
While recurrent networks can model sequential data, they suffer from key limitations:

\begin{itemize}
    \item difficulty capturing long-range dependencies,
    \item sequential computation limits parallelism,
    \item information is compressed into a single hidden state.
\end{itemize}

\medskip

\noindent
\textbf{Key question.}  
Can we design models that directly relate all elements of a sequence?

\medskip

\noindent
\textbf{Guiding idea.}  
Attention mechanisms allow each element of a sequence to interact with all others.

\subsection{Basic Attention Mechanism}

Attention computes a weighted combination of values.

\medskip

\noindent
Given:
\begin{itemize}
    \item queries $q$,
    \item keys $k_i$,
    \item values $v_i$,
\end{itemize}

define attention weights:
\[
\alpha_i = \frac{\exp(q \cdot k_i)}{\sum_j \exp(q \cdot k_j)}.
\]

\medskip

\noindent
Output:
\[
\text{Attention}(q, K, V) = \sum_i \alpha_i v_i.
\]

\medskip

\noindent
\textbf{Interpretation.}

\begin{itemize}
    \item keys determine relevance
    \item values carry information
    \item query selects what to focus on
\end{itemize}

\medskip

\noindent
\textbf{Insight.}  
Attention dynamically weights information based on relevance.

\subsection{Self-Attention}

In self-attention, queries, keys, and values are derived from the same sequence.

\medskip

\noindent
Given inputs $\{x_1, \dots, x_T\}$:

\[
q_i = W_Q x_i, \quad k_i = W_K x_i, \quad v_i = W_V x_i.
\]

\medskip

\noindent
Each position attends to all others:
\[
y_i = \sum_{j=1}^{T} \alpha_{ij} v_j.
\]

\medskip

\noindent
\textbf{Key properties.}

\begin{itemize}
    \item captures global dependencies
    \item fully parallelizable
\end{itemize}

\medskip

\noindent
\textbf{Insight.}  
Self-attention replaces sequential processing with global interaction.

\subsection{Scaled Dot-Product Attention}

To improve stability, attention uses scaling:

\[
\alpha_{ij} =
\frac{\exp\left(\frac{q_i \cdot k_j}{\sqrt{d}}\right)}
{\sum_{k} \exp\left(\frac{q_i \cdot k_k}{\sqrt{d}}\right)}.
\]

\medskip

\noindent
\textbf{Reason.}

\begin{itemize}
    \item prevents large dot products
    \item stabilizes gradients
\end{itemize}

\medskip

\noindent
\textbf{Insight.}  
Scaling ensures numerically stable similarity computation.

\subsection{Multi-Head Attention}

Instead of a single attention mechanism, we use multiple heads:

\[
\text{MultiHead}(Q,K,V) =
\text{Concat}(\text{head}_1, \dots, \text{head}_h) W_O.
\]

\medskip

\noindent
\textbf{Interpretation.}

\begin{itemize}
    \item each head learns different relationships
    \item captures multiple types of dependencies
\end{itemize}

\medskip

\noindent
\textbf{Insight.}  
Multi-head attention enriches representation by combining diverse interactions.

\subsection{Comparison with CNNs and RNNs}

\textbf{CNNs:}
\begin{itemize}
    \item local receptive fields
    \item limited long-range interaction
\end{itemize}

\textbf{RNNs:}
\begin{itemize}
    \item sequential processing
    \item information bottleneck
\end{itemize}

\textbf{Attention:}
\begin{itemize}
    \item global receptive field
    \item direct interactions between all positions
\end{itemize}

\medskip

\noindent
\textbf{Insight.}  
Attention removes structural constraints of locality and recurrence.

\subsection{Computational Considerations}

\textbf{Complexity.}

\begin{itemize}
    \item attention requires $O(T^2)$ operations
\end{itemize}

\medskip

\noindent
\textbf{Tradeoffs.}

\begin{itemize}
    \item more expressive interactions
    \item higher computational cost
\end{itemize}

\medskip

\noindent
\textbf{Insight.}  
Expressivity comes at a computational cost.

\subsection{Geometric Interpretation of Attention}

Attention can be interpreted geometrically as measuring similarity in a feature space.

\medskip

\noindent
Recall that attention weights are computed via dot products:
\[
\alpha_{ij} \propto \exp\left( \frac{q_i \cdot k_j}{\sqrt{d}} \right).
\]

\medskip

\noindent
\textbf{Geometric view.}

\begin{itemize}
    \item $q_i$ and $k_j$ are vectors in $\mathbb{R}^d$
    \item $q_i \cdot k_j$ measures alignment (similarity)
\end{itemize}

\medskip

\noindent
\textbf{Normalized interpretation.}

If vectors are normalized:
\[
q_i \cdot k_j = \|q_i\|\|k_j\|\cos\theta_{ij},
\]
so attention depends on the angle between vectors.

\medskip

\noindent
\textbf{Interpretation.}

\begin{itemize}
    \item large weight $\Rightarrow$ vectors aligned
    \item small weight $\Rightarrow$ vectors orthogonal or opposed
\end{itemize}

\medskip

\noindent
\textbf{Insight.}  
Attention selects information based on geometric similarity in representation space.

\subsection{Matrix Formulation of Attention}

Let:
\[
Q = [q_1, \dots, q_T]^\top, \quad
K = [k_1, \dots, k_T]^\top, \quad
V = [v_1, \dots, v_T]^\top.
\]

\medskip

\noindent
\textbf{Attention matrix.}

\[
A = \text{softmax}\left( \frac{QK^\top}{\sqrt{d}} \right).
\]

\medskip

\noindent
\textbf{Output.}

\[
Y = A V.
\]

\medskip

\noindent
\textbf{Interpretation.}

\begin{itemize}
    \item $QK^\top$ computes all pairwise similarities
    \item softmax produces a row-stochastic matrix
    \item each row of $A$ defines a weighted average over values
\end{itemize}

\medskip

\noindent
\textbf{Key observation.}

\begin{itemize}
    \item attention is a data-dependent linear operator
    \item weights depend on the input itself
\end{itemize}

\medskip

\noindent
\textbf{Insight.}  
Self-attention performs adaptive, input-dependent linear transformations.

\subsection{Attention as a Kernel Operator}

The attention mechanism resembles kernel methods.

\medskip

\noindent
Define a similarity function:
\[
k(x_i, x_j) = \exp\left( \frac{q_i \cdot k_j}{\sqrt{d}} \right).
\]

\medskip

\noindent
\textbf{Interpretation.}

\begin{itemize}
    \item attention computes a normalized kernel
    \item similar to kernel regression or smoothing
\end{itemize}

\medskip

\noindent
\textbf{Connection.}

\[
y_i = \sum_j \alpha_{ij} v_j
\]
resembles:
\[
f(x_i) = \sum_j k(x_i, x_j) y_j.
\]

\medskip

\noindent
\textbf{Key difference.}

\begin{itemize}
    \item classical kernels are fixed
    \item attention kernels are learned and data-dependent
\end{itemize}

\medskip

\noindent
\textbf{Insight.}  
Attention generalizes kernel methods by learning the similarity function.

\subsection{Continuous / Integral Operator View}

In the limit of large sequences, attention can be viewed as an integral operator.

\medskip

\noindent
Let $x(t)$ be a continuous input and define:
\[
y(t) = \int k(t,s) v(s) \, ds.
\]

\medskip

\noindent
\textbf{Interpretation.}

\begin{itemize}
    \item $k(t,s)$ measures interaction between positions $t$ and $s$
    \item $y(t)$ aggregates information across the domain
\end{itemize}

\medskip

\noindent
\textbf{Connection.}

\begin{itemize}
    \item discrete attention $\rightarrow$ finite-dimensional approximation
    \item continuous limit $\rightarrow$ operator on function spaces
\end{itemize}

\medskip

\noindent
\textbf{Relation to functional analysis.}

\begin{itemize}
    \item attention acts as a linear integral operator
    \item kernel defines the operator
\end{itemize}

\medskip

\noindent
\textbf{Insight.}  
Attention can be viewed as a learned operator acting on functions.

\subsection{Unifying Interpretation}

From these perspectives, attention can be understood as:

\begin{itemize}
    \item \textbf{Geometric:} similarity in feature space
    \item \textbf{Algebraic:} matrix-based linear transformation
    \item \textbf{Kernel-based:} adaptive similarity operator
    \item \textbf{Analytical:} integral operator on functions
\end{itemize}

\medskip

\noindent
\textbf{Core idea.}  
Attention is a flexible mechanism for modeling interactions across inputs.

\medskip

\noindent
\textbf{Key insight.}  
Modern architectures can be interpreted through multiple mathematical lenses, revealing deeper structure behind empirical success.

\subsection{A Unifying Perspective}

Attention mechanisms provide a new paradigm:

\begin{itemize}
    \item representations are computed through interactions,
    \item relationships between elements are learned dynamically.
\end{itemize}

\medskip

\noindent
\textbf{Core idea.}  
Learning involves selecting and combining relevant information across the entire input.

\medskip

\noindent
\textbf{Key insight.}  
Attention transforms representation learning into interaction modeling.

\subsection{Outlook}

In this lecture, we introduced:

\begin{itemize}
    \item attention mechanisms,
    \item self-attention,
    \item scaled dot-product attention,
    \item multi-head attention.
\end{itemize}

\medskip

\noindent
These ideas form the foundation of transformer architectures.

\medskip

\noindent
In the next lecture, we study transformers, which build entirely on attention mechanisms and enable large-scale models.

\medskip

\noindent
\textbf{Guiding principle.}  
Modeling interactions directly can overcome limitations of traditional architectures.